\documentclass[12pt,letterpaper]{article}
\usepackage[utf8]{inputenc}
\usepackage[spanish]{babel}
\usepackage{amsmath}
\usepackage{amsfonts}
\usepackage{amssymb}
\usepackage[margin=2.5cm]{geometry}

\begin{document}

\begin{center}
    \Large \textbf{Evaluación de Examen 2: Unidad II} \\
    \large Física para Ciencias de la Salud (005-1713) \\
    \vspace{0.5cm}
    \textbf{Estudiante:} Karla Rodríguez \quad \textbf{C.I.:} 33.628.014
    \vspace{0.3cm} \\
    \large \textbf{Calificación Final: 10/10 puntos}
\end{center}

\vspace{0.5cm}

\section*{Análisis de Resultados}

\subsection*{1. Cinemática (Aceleración media)}
\textbf{Procedimiento y resultado:} Extrae los datos con gran rigor formal al utilizar la notación delta para el tiempo ($\Delta t = 0,15 \text{ s}$). Plantea y sustituye correctamente la ecuación de aceleración, obteniendo el resultado analítico exacto de $4 \text{ m/s}^2$ con un correcto manejo del análisis dimensional. \\
\textbf{Veredicto:} \textbf{Correcto} (2/2 pts).

\subsection*{2. Dinámica (Leyes de Newton)}
\textbf{Procedimiento y resultado:} Muestra un despeje impecable a partir de la Segunda Ley de Newton ($a = \frac{F}{m}$). Sustituye los valores y divide $150 \text{ N}$ entre $25 \text{ kg}$, logrando el resultado perfecto de $6 \text{ m/s}^2$ y redactando una excelente conclusión final. \\
\textbf{Veredicto:} \textbf{Correcto} (2/2 pts).

\subsection*{3. Trabajo Mecánico}
\textbf{Procedimiento y resultado:} Destaca su nivel analítico al indicar de forma explícita que tanto la fuerza como el desplazamiento van ``hacia arriba'', justificando perfectamente que el ángulo sea $0^\circ$. Simplifica la fórmula trigonométrica y realiza el cálculo de $400 \text{ N} \cdot 0,8 \text{ m}$, obteniendo $320 \text{ J}$. \\
\textbf{Veredicto:} \textbf{Correcto} (2/2 pts).

\subsection*{4. Energía Potencial}
\textbf{Procedimiento y resultado:} Extrae de manera estructurada los datos del enunciado y aplica a la perfección la fórmula de la energía potencial gravitatoria ($E_p = m \cdot g \cdot h$). Es destacable que desarrolla un paso intermedio calculando el peso en Newtons ($11,76 \text{ N}$) antes de multiplicarlo por la altura ($1,5 \text{ m}$). El resultado de $17,64 \text{ J}$ es perfectamente exacto. \\
\textbf{Veredicto:} \textbf{Correcto} (2/2 pts).

\subsection*{5. Potencia Mecánica}
\textbf{Procedimiento y resultado:} Plantea correctamente la fórmula relacionando el trabajo entre la variación del tiempo ($P = \frac{W}{\Delta t}$). Ejecuta la sustitución de $75 \text{ J} / 60 \text{ s}$ logrando el resultado aritmético perfecto de $1,25 \text{ W}$. \\
\textbf{Veredicto:} \textbf{Correcto} (2/2 pts).

\vspace{0.5cm}
\section*{Comentarios Generales y Sugerencias}
¡Felicidades por obtener la calificación máxima! El examen evidencia un dominio total de los conceptos teóricos y una excelente destreza matemática para plasmarlos en el papel.
\begin{itemize}
    \item El formato de resolución utilizado (separando de forma impecable en \textit{Datos}, \textit{Fórmula}, desarrollo y conclusión redactada) es el estándar oro esperado en reportes de laboratorio e investigación clínica.
    \item Se valora enormemente el \textbf{nivel de rigor teórico empleado}, como el uso del símbolo $\Delta t$ para el intervalo de tiempo en las preguntas 1 y 5, así como el análisis direccional vectorial en la pregunta 3.
    \item Como única y pequeña observación formal de cara a futuras publicaciones: recordar que en el Sistema Internacional (S.I.) el símbolo para el prefijo \textit{kilo} siempre se escribe en minúscula ($\text{kg}$ en lugar de $\text{Kg}$).
    \item ¡Continúa con este estupendo nivel de excelencia!
\end{itemize}

\end{document}